\section{Monte Carlo event generators}
\label{sec:generators}

In this section we present the MC event generators that will be considered in this benchmark.
%
We focus on three NLO-accurate general-purpose MC generators, namely \sherpa, \herwig, and \powheg. 
%
For inclusive observables and charm production, their predictions are compared with the analytical structure-function calculations provided by \yadism.
%
For completeness, we have also studied predictions provided by the leading-order (LO) MCs \pythia and \mgamc, finding excellent agreement with the other MCs when run at LO.\footnote{Results for \pythia and \mgamc are not shown in this work but are available from the project repository, see App.~\ref{app:implementation}.}
%
Furthermore, we consider event generators tailored to neutrino physics, with emphasis on \genie, the current default MC in FASER and SND@LHC, as well as {\sc\small NuWro} and {\sc\small GiBUU} (the latter two in App.~\ref{app:nugen}).
%
Table~\ref{tab:generators} indicates the main features of these generators.
%
For each generator we indicate the version used, the source of the hard process, the QCD perturbative order, and the shower and hadronisation algorithms.
%
We now provide more details on each of these generators in turn.

%%%%%%%%%%%%%%%%%%%%%%%%%%%%%%%%%%%%%%%%%%
\begin{table}[!tbp]
  \centering
  \small
  \renewcommand{\arraystretch}{1.60}
  \setlength{\tabcolsep}{4pt}
  \resizebox{\textwidth}{!}{%
  \begin{tabular}{llllll}
    \toprule
    Generator & Version & Hard process & Order & Shower & Hadronisation \\
    \midrule
    \pythia      & 8.311  & internal DIS ME        & LO            & $p_T$-ordered, dipole recoil & Lund string \\
                 &        &                        &               & (also {\sc Vincia}~\cite{Fischer:2016vfv}, {\sc Dire}~\cite{Hoche:2015sya}) & \\
    \mgamc       & 3.7.2  & tree-level ME          & LO            & {\sc\small Pythia}~8.311 & Lund string \\
                 &        &                        &               & (also parton level) & \\
    \genie       & 3.06.02 & {\sc EMDIS} / {\sc DIS-CC} (+charm)  & LO (Bodek--Yang) & none & {\sc AGKY} + {\sc\small Pythia}~8.312 \\
                 &        & tune {\tt G18\_02a}    &               &  & \\
    \midrule 
    \sherpa      & 3.0.5  & {\sc Comix} / {\sc Amegic}~\cite{Krauss:2001iv} & NLO ({\sc MC@NLO}) & CS dipole (also {\sc Dire}) & {\sc AHADIC} cluster \\
    \herwig      & 7.3.0  & {\sc MEDISNC} / {\sc MEDISCC} & NLO ({\sc POWHEG}) & angular-ordered & cluster \\
    \powhegres   & {\tt DIS\_v} & POWHEG-BOX-RES   & NLO           & {\sc\small Pythia}~8.311 & Lund string \\
                 &        &                        &               & or {\sc\small Herwig}~7.3.0 & or cluster \\
    \powhegtwo   & {\tt nu-DIS} & POWHEG-BOX-V2    & NLO           & {\sc\small Pythia}~8.311 & Lund string \\
                 &        &                        &               & or {\sc\small Herwig}~7.3.0 & or cluster \\
    \genie       & 3.06.02 & {\sc HEDIS}, tune {\tt GHE19\_00a} & NLO (\apfel)  & none & {\sc LEPTO}-type + {\sc\small Pythia}~8.312 \\
    \midrule
    \yadism      & 0.13.11 & DIS structure functions & LO, NLO, NNLO & --- & --- \\
    \bottomrule
  \end{tabular}}
  \vspace{0.5cm}
  \caption{LO and NLO event generators considered in this benchmark.
  %
    For each generator we indicate the version, the source of the hard process (ME: matrix element), the QCD perturbative order, and the shower and hadronisation algorithms (CS: Catani--Seymour).
  %
  \sherpa and \herwig can also be run at LO through the same matrix elements.
  %
  In this work both \genie configurations perform hadronisation and decays with \pythia rather than with {\sc Pythia6}, the \genie default.
  %
  The last row corresponds to the analytical structure-function calculation of \yadism, interfaced to {\sc EKO}~0.15.5~\cite{Candido:2022tld} for the DGLAP evolution.
  %
  Appendix~\ref{app:nugen}
  presents results for two other neutrino generators,  {\sc\small NuWro} and {\sc\small GiBUU}.  
  }
  \label{tab:generators}
\end{table}
%%%%%%%%%%%%%%%%%%%%%%%%%%%%%%%%%%%%%%%%%%

\paragraph{\pythia.}  {\sc\small Pythia}~8.311~\cite{Bierlich:2022pfr} simulates DIS processes
through its internal LO neutral- and charged-current matrix elements, showered
with the $p_T$-ordered shower and hadronised with the Lund string model~\cite{Andersson:1983ia}.  
%
The baseline shower uses the dipole-recoil scheme~\cite{Cabouat:2017rzi}, in which initial-state radiation recoils against the colour dipole rather than the scattered lepton.
%
In this work, \pythia is considered both as a stand-alone generator and as the provider of shower and hadronisation algorithms for some of the NLO generators.
%
Unless otherwise specified, here \pythia is run with the Monash tune~\cite{Skands:2014pea} for non-perturbative physics. Other tunes, such as the forward physics tune of~\cite{Fieg:2023kld}, can be easily implemented. 

\paragraph{\sherpa.}  \sherpa~v3.0.5~\cite{Sherpa:2019gpd,Bothmann:2024} provides either an LO sample from {\sc Comix}~\cite{Gleisberg:2008fv} or an NLO sample using the MC@NLO~\cite{Frixione:2002ik} scheme for the matching to the Catani--Seymour (CS) dipole
shower~\cite{Schumann:2007mg}.
%
Hadronisation in \sherpa is based on the {\sc AHADIC} cluster model~\cite{Winter:2003tt,Chahal:2022rid}.  

\paragraph{\herwig.}  {\sc\small Herwig}~7.3.0~\cite{Bellm:2015jjp,Bewick:2023tfi} provides DIS simulations at LO through the {\sc MEDISNC} and {\sc MEDISCC} modules, and at NLO through its POWHEG implementation matched to the angular-ordered shower, and uses a cluster hadronisation model~\cite{Webber:1983if}.
%
\herwig admits two independent NLO implementations.  The first matches {\sc
Matchbox}~\cite{Platzer:2011bc} one-loop amplitudes to the shower in the MC@NLO scheme; the second is
\herwig's native POWHEG implementation, built on the same matrix-element classes as
the LO calculation, and is the one adopted here. 
%
The value of the intrinsic $k_T$ used in the \herwig simulations is adjusted to match the DIS kinematics. 

\paragraph{\powheg.}  Two independent POWHEG~\cite{Nason:2004rx,Frixione:2007vw,Alioli:2010xd}
implementations are used in this work.
%
First,
\powhegres is the
{\tt DIS\_v} process~\cite{Banfi:2023mhz} of POWHEG-BOX-RES~\cite{Jezo:2015aia}, covering inclusive and charm
production in both currents with mass effects neglected; it supplies the NLO
entry for NC muon scattering. 
%
Second, \powhegtwo is the {\tt nu-DIS}
process of POWHEG-BOX-V2~\cite{Buonocore:2024pdv}, developed for neutrino DIS with a massive
final-state charm quark, and supplies the NLO entry for charged-current
neutrino scattering.  
%
By default both are showered and hadronised with \pythia, and hence with
Lund-string hadronisation; showering and hadronisation with \herwig is studied as a variant.
%
Mass effects in \powhegtwo are only available for the generation of charm-tagged events, while for inclusive DIS production charm-quark mass effects are neglected. 
%
For inclusive generation, both \powheg variants agree with each other and with the ZM-VFNS \yadism calculation at the per-cent level at NLO. 
%
See also~\cite{Castro:2026xyr} for a recent POWHEG-BOX-RES calculation of heavy-quark pair production in NC DIS.

\paragraph{\mgamc.}  \mgamc~3.7.2~\cite{Alwall:2014hca} is used to generate events at LO matched to \pythia.
%
In its current implementation, this generator cannot be used to simulate DIS events at NLO accuracy.

\paragraph{\genie.}   
%
As opposed to the other generators considered here, the \genie~\cite{Andreopoulos:2009rq,Andreopoulos:2015wxa} settings (such as PDF choice or matrix-element (ME) model) cannot be straightforwardly modified and are only available for a fixed set of tunes. 
%
Here LO simulations for neutrino scattering are produced with \genie using 
the {\tt G18\_02a\_00\_000} tune~\cite{GENIE:2021zuu}, which is the same configuration that FASER uses~\cite{FASER:2024ykc}.
%
This tune is built on the Bodek--Yang model~\cite{Bodek:2002vp,Bodek:2003wd} with
GRV98LO~\cite{Gluck:1998xa}.  
%
For NC muon scattering we use the
{\sc EMDIS} channel of that tune.
%
For CC scattering we also consider the {\sc HEDIS} module~\cite{Soto:2021vdc} of {\tt GHE19\_00a}, built on BGR18 structure functions~\cite{Bertone:2018dse} at NLO through \apfel~\cite{Bertone:2013vaa}, achieving NLO accuracy for inclusive distributions.
%
This tune evaluates its structure functions with the {\tt NNPDF31sx\_nlo\_as\_0118\_LHCb\_nf\_6} set rather than with \nnpdf.
% 
To separate PDF effects from other model parameters, we also generate LO events where GRV98
is replaced by \nnpdf.

\genie does not simulate a parton shower, either initial- or final-state, for the configurations considered here: the struck quark and the colour-connected target remnant are passed directly to string fragmentation, so \genie final states contain no perturbative QCD radiation beyond the LO partonic process.
%
In the default tune the hadronic system is handled by the {\sc AGKY} model~\cite{Yang:2009zx}, which takes hadron multiplicities and momenta from the KNO parametrisation~\cite{Koba:1972ng} below a hadronic invariant mass of $2.3$~GeV and from \pythia string fragmentation above $3$~GeV, interpolating in between; {\sc HEDIS} instead uses a {\sc LEPTO}-derived~\cite{Ingelman:1996mq} interface to \pythia at all $W$.
%
In {\sc HEDIS} the NLO corrections enter only as a reweighting of the LO cross-section, so its final states are LO in structure. 
%
Furthermore, \genie treats charm production\footnote{In the course of this work we found that the charm-production DIS cross-section of \genie multiplies by the $W$-boson propagator factor $(1+Q^2/M_W^2)$ instead of dividing by it. 
%
This overestimates the charm cross-section for $\nu_\mu$ scattering on a proton by about $4\%$ at $E_\nu=400$~GeV and $8\%$ at 1~TeV, growing rapidly at higher energies, while the fiducial inclusive cross-section changes by less than $0.2\%$. All \genie predictions in this work are obtained with the corrected cross-section.}  as a separate process~\cite{Aivazis:1993pi} and subtracts it from the inclusive DIS cross-section internally.  
%
The charm channel must therefore always be generated alongside the
inclusive one.
%
Concerning NC scattering, the {\sc EMDIS} channel
list is a three-flavour calculation where charm is absent both from the initial and final state. 

An advantage of \genie as compared to the general-purpose generators is that it can account also for non-DIS processes such as quasi-elastic, resonance, and diffractive scattering which are then combined with the DIS contribution.
%
It can also cover the shallow-inelastic scattering (SIS) region at low $Q^2$, which provides a non-negligible contribution to total event yields~\cite{Candido:2023utz}.
%
In App.~\ref{app:nondis} we quantify the importance of non-DIS contributions for the FASER neutrino kinematics, finding that these effects can in general be neglected for all relevant energies provided one restricts the calculation to the DIS region {\it e.g.} by applying $Q^2>4$~GeV$^2$ and $W>3$~GeV cuts. 
%
Whenever necessary, predictions simulated with NLO generators can be complemented with the contributions from outside the DIS region (the SIS and resonance regions) as estimated by \genie using the pipeline presented in this work.
%
This can be done, for instance, by adding the \genie prediction for the kinematic region excluded by the DIS cuts, evaluated for the same selection, with a smooth interpolation across the boundary where differential distributions require it.
%
An example of this is presented in Sect.~\ref{sec:faser-app} for the comparison with the FASER electronic-detector muon-neutrino data, where the $Q^2<4$~GeV$^2$ contribution from the SIS region is taken from \genie.

\paragraph{\yadism.} Analytic calculations of the DIS structure functions up to NNLO and approximate N$^3$LO (aN$^3$LO)~\cite{NNPDF:2024nan} and accounting for heavy-quark mass effects are provided by
\yadism~\cite{Candido:2024rkr}, using the same PDFs, renormalisation and factorisation scales and electroweak inputs as the MC generators. 
%
Cross-sections are assembled from $F_2$, $F_L$ and $xF_3$ and then integrated over the
fiducial region.
%
\yadism further allows the heavy-quark scheme to be varied, comparing massless calculations in the zero-mass variable-flavour-number scheme (ZM-VFNS) with the FONLL-matched ones~\cite{Forte:2010ta} which account for heavy-quark mass effects.
%
\yadism has been benchmarked with other widely used DIS tools such as \apfel~\cite{Bertone:2013vaa} and {\sc\small HOPPET}~\cite{Salam:2008qg,Karlberg:2025hxk}.